%% =====================================================================
%%  T-06-04  Behaviour Over Declaration
%%  -------------------------------------------------------------------
%%  One idea per file. Core is mandatory. Slots in canonical order.
%%  Max 700 words. No layout commands. No filler. New source material
%%  for this entry goes in sources/<BOOK>/T-06-04.tex -- never by
%%  rewriting this file (docs/RULES.md R0.1).
%% =====================================================================
%% @kind: topic
%% @id: T-06-04
%% @title: Behaviour Over Declaration
%% @subtitle: What a person does is the only statement they cannot edit
%% @chapter: c06
%% @order: 40
%% @status: stable
%% @sources: B1, B3
%% @updated: 2026-09-19

\topic{T-06-04}{Behaviour Over Declaration}{What a person does is the only statement they cannot edit}

\begin{core}
Words are chosen for effect and can be revised at no cost. Conduct is expensive to fake over time. Where the two disagree, the disagreement is the information --- and the declaration is the less reliable of the pair.
\end{core}
\begin{mechanism}
Declarations are cheap because they are produced for an audience and cost the speaker nothing. Conduct is expensive because it is repeated, observable by many people over time, and inconsistent with hidden intentions in ways that eventually show. A person can maintain a stated position indefinitely; maintaining a contradictory pattern of action is much harder. This is why the reliable reading is always the longitudinal one --- behaviour sampled across situations and across time --- and why a single conversation, however revealing it feels, is weak evidence. It is also why demonstration beats argument: an action that changes someone's position without a word leaves no grievance and no counter-argument to work against.
\end{mechanism}
\begin{tells}
  \item A person's stated values predict nothing about their actual choices.
  \item Commitments are made fluently and met loosely.
  \item Their account of their own conduct is consistently more generous than other people's accounts of it.
  \item You find yourself persuaded by what they said rather than by what happened.
  \item The gap widens under pressure rather than closing.
\end{tells}
\begin{moves}
  \item Sample behaviour across at least three different situations before forming a view, including one where the person has something to lose.
  \item Where you need agreement, produce a result rather than an argument.
  \item Say less than the situation invites, and let your actions carry the position.
  \item Where you must state a position, state it once and then demonstrate it rather than repeating it.
\end{moves}
\begin{counters}
  \item Weigh the last six months of what someone did above everything they said this week.
  \item Ask for a small, cheap, immediate action that follows from their declaration. Refusal or delay on something trivial is more diagnostic than any explanation.
  \item Track commitments in writing. The pattern emerges by itself and does not need to be argued about.
  \item Do not accept a revised account of past conduct as evidence about future conduct.
\end{counters}
\begin{cost}
Behaviour is also interpretable, and interpretations differ --- a person may be unreliable rather than dishonest, constrained rather than unwilling. Reading conduct requires context that outsiders often lack, and the confident behaviourist is frequently just as wrong as the credulous listener, only later. The rule also breaks down for one-off interactions, which have no longitudinal record to read; there, structural protections matter more than any judgement of character.
\end{cost}

\sources{B1, B3}
\seesources
\seealso{T-06-01, T-06-03, T-16-01, T-23-01}
